\chapter{Sviluppi Futuri}

Il lavoro presentato in questa fase costituisce le fondamenta formali
e l'infrastruttura architetturale del sistema di analisi
\textit{language-agnostic}. Le fasi di estrazione, risoluzione e
costruzione del grafo hanno dimostrato la validità dell'approccio
basato su tipi e composizione funzionale. Tuttavia, per raggiungere
un grado di maturità tale da poter analizzare codebase complesse e
reali in produzione, il sistema richiederà diverse estensioni nei
prossimi mesi di sviluppo.

Di seguito vengono delineati i principali filoni di ricerca e
implementazione futuri.

\section{Estensione del Supporto ai Costrutti del Linguaggio}
Attualmente, la Rappresentazione Intermedia (IR) copre i costrutti
fondamentali della programmazione orientata agli oggetti e
procedurale (Classi, Interfacce, Funzioni, Metodi). L'IR dovrà essere
espansa per supportare costrutti più avanzati che impattano in modo
significativo sull'architettura:
\begin{itemize}
  \item \textbf{Generics e Polimorfismo Parametrico}: Linguaggi come
    Java, C++ (Template) e Rust fanno un uso massiccio di tipi
    generici. Il sistema dei tipi $\mathcal{T}$ dovrà essere esteso
    per catturare i parametri di tipo (es. \texttt{List<T>}) e i
    vincoli associati (es. \texttt{T: Clone} in Rust o \texttt{T
    extends Serializable} in Java).
  \item \textbf{Tipi Somma Alggebrici}: L'introduzione di
    \textbf{Enums} complessi (con dati associati, tipici di Rust o
    Swift) e \textbf{Unions} (in C/C++) richiederà la definizione di
    un nuovo Componente nell'IR, distinto dalle normali Classi o Struct.
\end{itemize}

\section{Supporto a Paradigmi Architetturali Avanzati}
Molti linguaggi moderni utilizzano costrutti che alterano la
struttura o il comportamento del codice in fase di compilazione o a
runtime. Ignorare questi elementi significherebbe perdere dipendenze
architetturali cruciali:
\begin{itemize}
  \item \textbf{Decorators e Annotazioni}: In Python e Java, i
    decoratori (\texttt{@Entity}, \texttt{@Autowired}) spesso
    iniettano dipendenze silenziose (es. verso framework come Spring
    o Django). Sarà necessario estrarre questi metadati e mapparli
    come speciali archi di dipendenza.
  \item \textbf{Mixins e Traits con Implementazioni di Default}:
    Sebbene i Trait siano abbozzati, la gestione dei Mixins (comuni
    in linguaggi dinamici o in Dart) richiede una logica di
    risoluzione più sofisticata per tracciare correttamente
    l'ereditarietà multipla o l'iniezione di comportamento.
  \item \textbf{Programmazione Orientata agli Aspetti (AOP)}:
    Costrutti come gli \textit{Aspects} in Java modificano il flusso
    di esecuzione ortogonalmente alla gerarchia delle classi.
    Tracciarli nel Grafo delle Dipendenze rappresenterà una sfida
    interessante per il mantenimento dell'agnosticismo.
\end{itemize}

\section{Affinamento delle Euristiche di Estrazione}
L'approccio attuale si basa su funzioni di dispatch che traducono i
nodi del Concrete Syntax Tree (CST) di Tree-sitter. Nei prossimi
mesi, le euristiche dovranno essere irrobustite per gestire:
\begin{itemize}
  \item \textbf{Costrutti Sintattici Ambigui}: Migliorare il parsing
    di macro (C/C++ e Rust) che espandendosi generano nuovo codice
    non visibile nell'AST originale.
  \item \textbf{Graceful Degradation}: Implementare meccanismi più
    solidi per estrarre informazioni parziali quando il parser
    incontra porzioni di codice gravemente malformate o non
    supportate, garantendo che l'analisi globale non fallisca ma
    segnali puntualmente il degrado.
\end{itemize}

\section{Arricchimento del Grafo e Analisi delle Metriche}
Una volta che l'astrazione sarà sufficientemente ricca, l'ultimo step
riguarderà lo sfruttamento del Grafo delle Dipendenze ($\mathcal{G}$):
\begin{itemize}
  \item \textbf{Nuove Semantiche di Dipendenza}: Espandere l'insieme
    $\text{Kind}$ degli archi per includere relazioni temporali o di
    ownership (es. il passaggio di referenze mutabili in Rust vs il
    passaggio per valore in C).
  \item \textbf{Metriche Architetturali}: Sviluppare algoritmi che,
    navigando il grafo, calcolino metriche classiche (Accoppiamento,
    Coesione, Cicli di Dipendenza) in modo totalmente indipendente
    dal linguaggio sorgente analizzato.
\end{itemize}
